\documentclass[10pt]{beamer}
\usetheme{metropolis} % A modern, clean theme
\usepackage{booktabs}
\usepackage{pgfplots}
\pgfplotsset{compat=1.18}

\title{Experiment Results: Kulishova-Replication}
\subtitle{Machine Learning vs. Deep Learning Benchmarking}
\date{\today}
\author{Data Export Analysis}
\institute{W&B Project: Binary Classification}

\begin{document}

\maketitle

\begin{frame}{Project Overview}
    \begin{itemize}
        \item \textbf{Dataset:} Kulishova Extracted Features.
        \item \textbf{Objective:} Compare classical ML models with Deep Learning for binary classification.
        \item \textbf{Pre-processing:} Majority class downsampling to address class imbalance.
        \item \textbf{Models Tested:} 
            \begin{itemize}
                \item Support Vector Machine (SVM)
                \item Random Forest (RF)
                \item 1D-Convolutional Neural Network (1D-CNN)
            \end{itemize}
    \end{itemize}
\end{frame}

\begin{frame}{Model Performance Comparison}
    The 1D-CNN outperformed classical methods, though all models showed high stability.
    \vspace{0.5cm}
    \begin{table}
        \centering
        \begin{tabular}{lccc}
            \toprule
            \textbf{Metric} & \textbf{1D-CNN} & \textbf{Random Forest} & \textbf{SVM} \\
            \midrule
            Accuracy (Val) & 94.04\% & 92.87\% & 91.36\% \\
            F1-Score (Val) & 93.64\% & 92.86\% & 91.36\% \\
            Best Epoch & 17 & N/A & N/A \\
            \bottomrule
        \end{tabular}
        \caption{Validation metrics across different architectures.}
    \end{table}
\end{frame}

\begin{frame}{Performance Visualization}
    \begin{center}
    \begin{tikzpicture}
        \begin{axis}[
            ybar,
            enlarge x limits=0.5,
            legend style={at={(0.5,-0.15)}, anchor=north, legend columns=-1},
            ylabel={Score (\%)},
            symbolic x coords={1D-CNN, RF, SVM},
            xtick=data,
            nodes near coords,
            nodes near coords align={vertical},
            ylim={85, 100},
            height=6cm,
            width=9cm,
            bar width=15pt,
        ]
            \addplot coordinates {(1D-CNN, 94.04) (RF, 92.87) (SVM, 91.36)};
            \addplot coordinates {(1D-CNN, 93.64) (RF, 92.86) (SVM, 91.36)};
            \legend{Accuracy, F1-Score}
        \end{axis}
    \end{tikzpicture}
    \end{center}
\end{frame}

\begin{frame}{1D-CNN Training Characteristics}
    The deep learning model utilized a 1D-CNN architecture optimized for feature sequences.
    
    \begin{itemize}
        \item \textbf{Loss Function:} Binary Crossentropy.
        \item \textbf{Activation:} Sigmoid (Output).
        \item \textbf{Optimizer Progress:} 
            \begin{itemize}
                \item Final Training Loss: 0.1590
                \item Final Validation Loss: 0.1455
                \item Final Learning Rate: $\sim$0.1398
            \end{itemize}
        \item \textbf{Efficiency:} Peak performance reached at Epoch 17 of 20.
    \end{itemize}
\end{frame}

\begin{frame}{Generalization & Overfitting Analysis}
    \begin{block}{Random Forest}
        High consistency between Train F1 (0.9455) and Val F1 (0.9286), indicating excellent generalization.
    \end{block}
    
    \begin{block}{Support Vector Machine}
        Perfect parity between training and validation metrics (91.36\%), suggesting the model has captured the underlying linear or kernel-based variance without any overfitting.
    \end{block}
\end{frame}

\begin{frame}{Conclusion \& Recommendations}
    \begin{itemize}
        \item \textbf{Primary Choice:} 1D-CNN is recommended for maximum predictive accuracy (94.04\%).
        \item \textbf{Alternative:} Random Forest provides a competitive (92.87\%) and more interpretable alternative with lower computational overhead.
        \item \textbf{Future Steps:} Evaluate model robustness against non-downsampled (imbalanced) real-world data to verify recall on the minority class.
    \end{itemize}
\end{frame}

\end{document}